\section{Comparison to Existing Work}
\label{sec:existing-work}

% ============================================================
\subsection{GDS versus EPA-Based Frameworks}
\label{sec:gds-vs-epa}
% ============================================================

The most widely cited team quality metric in professional NFL analytics is
Football Outsiders' Defense-adjusted Value Over Average (DVOA). Like the \gds{}
framework developed in this paper (\Cref{sec:gds}), DVOA decomposes team
performance into offensive, defensive, and special teams components, enabling
analysts to identify which phase of the game drives a team's overall quality.
Both approaches share the fundamental insight that a single aggregate number
obscures the strategic composition of team strength, and that meaningful
evaluation requires examining the constituent parts. Qualitatively, the two
frameworks would identify the same teams as elite -- both would recognize the
2023 San Francisco 49ers as a balanced elite team and the 2024 Baltimore
Ravens as an offensively dominant contender -- and both would reach similar
conclusions regarding the primacy of offensive quality in championship
contention.

The differences, however, are substantial and methodologically motivated. DVOA
operates at the play level, assigning value to each individual snap based on
Expected Points Added (EPA). \gds{} operates at the drive level, using
expected points derived from a multinomial outcome model as its currency of
evaluation (\Cref{sec:xscore}). This
distinction is not merely one of granularity; it reflects a different
theoretical commitment about the fundamental unit of football scoring
opportunity. A single explosive play -- a broken coverage resulting in a
75-yard touchdown, or a fumble on the first snap -- can dramatically distort
play-level EPA totals for a game. Drive-level aggregation smooths this noise
by asking whether the overall sequence of plays produced a scoring outcome
that exceeded or fell short of expectations given the starting field position.
In this sense, \gds{} captures sustained execution quality rather than the
sum of isolated individual events, providing a measure more reflective of
systematic team capability.

A second critical distinction concerns transparency and reproducibility. DVOA
is proprietary \parencite{schatz2005dvoa}: neither the precise opponent adjustments, the weighting
schemes, nor the underlying data transformations are publicly documented in
sufficient detail to permit independent replication. Researchers citing DVOA
must treat it as a black box, accepting its outputs on trust. The \gds{}
framework, by contrast, is fully open-source. The \xscore{} model, the
isotonic calibration procedure, the \xvoa{} decomposition, and the archetype
classification are all implemented in Python using publicly available
\texttt{nflfastR} data \parencite{carl2021nflfastr}, and the complete codebase
accompanies this paper. Any researcher can replicate, extend, or critique
the results. In an era where reproducibility concerns pervade quantitative
social science, this transparency represents a meaningful contribution.

The primary advantage DVOA retains is opponent adjustment: each play's value
is scaled by the quality of the opposing defense or offense, creating a
measure that accounts for strength of schedule. \gds{} does not incorporate
such an adjustment (\Cref{sec:limitations}), meaning that a team accumulating
\xvoa{} against weak opponents is not penalized relative to one facing elite
competition. Over a 17-game regular season, strength-of-schedule effects
partially average out, but this remains a limitation relative to DVOA's
approach. Future extensions of the framework could incorporate opponent
adjustment while preserving the drive-level, open-source character that
distinguishes \gds{}.

% ============================================================
\subsection{Comparison to Robst, VanGilder, and Berri (2011)}
\label{sec:comparison-robst}
% ============================================================

\textcite{robst2011defense} represents the most direct academic predecessor to
the present investigation. Their study examined whether defensive metrics
predict NFL playoff success, concluding that offensive statistics are
substantially stronger predictors and that the ``defense wins championships''
maxim lacks empirical support. The directional conclusion of this paper is
entirely consistent with their finding: across all five statistical procedures
employed in \Cref{ch:results}, offensive quality dominates defensive quality as
a predictor of postseason success.

The methodological advance offered by \gds{} over Robst et al.'s approach is
considerable, however. Their analysis relied on traditional box-score statistics
 -- total yards gained, total yards allowed, points scored, points allowed -- that
conflate quality with volume and fail to account for game context. A team
protecting a large lead in the fourth quarter will frequently allow ``garbage
time'' yards and points that inflate its defensive statistics without reflecting
genuine defensive weakness. Similarly, a team that plays with a lead controls
possession through running, suppressing its own passing yardage totals without
implying offensive inadequacy. The \xscore{} model addresses this directly:
by conditioning drive-outcome probabilities on field position, down, distance,
score differential, and time remaining (\Cref{sec:xscore}), the \xvoa{} metric
captures performance relative to a calibrated baseline that accounts for game
state. A team allowing a late touchdown when trailing by four scores receives
minimal defensive penalty, because the expected scoring probability was already
elevated by the situation. This context-sensitivity represents a substantial
refinement over raw box-score approaches and provides a more rigorous
measurement basis for the offensive primacy conclusion.

% ============================================================
\subsection{Novel Contributions}
\label{sec:novel-contributions}
% ============================================================

Four contributions distinguish this paper from existing work on NFL team
evaluation and the ``defense wins championships'' debate.

First, \gds{} is, to the author's knowledge, the first fully open-source,
reproducible team quality metric with an explicit three-component decomposition
into offensive, defensive, and special teams value. While
\textcite{yurko2019nflwar} developed the open-source \texttt{nflWAR} framework
for player-level EPA decomposition, their focus was individual attribution
rather than team-level quality measurement. The team-level metrics most commonly
used by practitioners -- DVOA, ESPN's QBR-based team ratings, and Pro Football
Focus grades -- are uniformly proprietary. By providing a transparent,
replicable alternative, this paper enables other researchers to build upon,
critique, or extend the framework without licensing barriers.

Second, the choice of drive-level probability modeling rather than play-level
EPA represents a theoretically motivated departure from the dominant paradigm
in football analytics. The drive is the fundamental unit of scoring opportunity
in American football: a team receives possession, executes a sequence of plays,
and either scores or does not. The \xscore{} model respects this structure by
predicting the probability of scoring a touchdown on each drive given its
initial conditions, then measuring whether the observed outcome exceeded or
fell short of that expectation. This approach avoids the well-documented issues
with play-level EPA, including sensitivity to individual plays, instability in
small samples, and the arbitrary allocation of value between consecutive plays
in a sequence \parencite{yurko2019nflwar}.

Third, this paper provides the first systematic, multi-method test of the
``defense wins championships'' hypothesis using a calibrated expected-value
framework. Previous investigations relied on either traditional statistics
\parencite{robst2011defense} or purely predictive models
\parencite{lock2014randomforests} without an explicit probability baseline.
The combination of 224 team-seasons across seven seasons (2018--2024), five
distinct statistical procedures (Spearman rank correlation, Pearson
correlation with $r^2$ decomposition, binary logistic regression, OLS
regression, and Cohen's $d$), and
pre-registered directional hypotheses provides considerably stronger evidence
than any single prior study. The convergent rejection across all methods
strengthens the inferential warrant beyond what any individual test could
provide.

Fourth, the ``competent defense threshold'' concept -- the finding that
defensive quality operates as a necessary but insufficient condition for
championship success, with diminishing returns beyond a moderate level -- has
not been previously formalized in the academic literature. This nuanced
conclusion moves beyond the binary framing of ``offense versus defense'' to
articulate a strategic principle: elite teams invest in defensive competence
to avoid catastrophic vulnerability, but allocate marginal resources toward
offensive excellence where the return on investment is substantially higher.

% ============================================================
\subsection{Consistency with Broader Literature}
\label{sec:consistency-literature}
% ============================================================

The offensive primacy finding is consistent with a growing body of evidence
across multiple methodological traditions. \textcite{lock2014randomforests}
employed random forests to predict NFL game outcomes and found that offensive
variables -- particularly passing efficiency metrics -- were consistently more
predictive of wins than defensive variables. Their machine learning approach,
operating at a different level of analysis and using a different modeling
strategy, reached the same directional conclusion as the present work,
suggesting that the finding is robust to methodological choices rather than
an artifact of any particular analytic framework.

\textcite{massey2013loserscurse} approached the question from a labor
economics perspective, demonstrating that NFL teams systematically overvalue
early draft picks used on defensive players. Their ``Loser's Curse'' finding
implies a market-level recognition that defensive talent is overpriced
relative to its contribution to winning -- a conclusion that aligns precisely
with the empirical finding that defensive \xvoa{} explains less than half the
variance in winning that offensive \xvoa{} does. If defensive quality truly
drove championships, one would expect the market to price defensive talent
at a premium reflecting that value; instead, the observed overvaluation
suggests that decision-makers are paying for a narrative (``defense wins
championships'') rather than a statistical reality.

More broadly, the conceptual framework of using expected scoring models to
evaluate team quality extends the approach pioneered by
\textcite{rathke2017expectedgoals} in association football. Rathke's expected
goals (xG) model demonstrated that comparing observed goals to
probability-based expectations reveals systematic over- and under-performance
attributable to team quality rather than luck. The \xscore{} model applies
this same logic to American football drives: observed touchdowns minus
expected touchdowns (i.e., \xvoa{}) isolates the quality signal from
situational noise. The successful application across sports demonstrates
the portability of probability-based performance measurement frameworks and
suggests that the drive-level approach developed here could be adapted to
other invasion sports where possession sequences culminate in discrete
scoring events.
